\section{Motor model}\label{sec:motor}

In \cref{sec:control} the controllers impose the torque directly. The robot
has two brushed DC motors with a gearbox, driven by a voltage $V$ through an
H-bridge; the torque they produce depends on $V$ and on the speed of the
wheel. This section derives the model used in the simulation.

\subsection{Datasheet values}

One motor, rated voltage \SI{12}{V}, gear ratio 21.3:1 (\cref{tab:motor}).

\begin{table}[h]
\centering
\caption{Datasheet values of one motor (output shaft).}
\label{tab:motor}
\begin{tabular}{@{}lll@{}}
\toprule
Quantity & Value & SI \\
\midrule
Rated voltage $V_{\mathrm{nom}}$ & \SI{12}{V} & \\
No-load speed $\omega_0$ & \SI{280}{rpm} & \SI{29.32}{rad/s} \\
No-load current & \SI{60}{mA} & \\
Stall torque $\tau_s$ & \SI{1.7}{kg.cm} & \SI{0.1667}{N.m} \\
Stall current $i_s$ & \SI{1.3}{A} & \\
\bottomrule
\end{tabular}
\end{table}

\subsection{Equations}

\paragraph{Electrical equation.} The winding has resistance $R_m$ and
inductance $L_m$. When the rotor turns, it generates a back electromotive
force proportional to its speed, which opposes the applied voltage. Referring
all quantities to the output shaft of the gearbox, with $\omega_{\mathrm{rel}}$
the speed of the output shaft relative to the motor housing, Kirchhoff's
voltage law for the winding gives
\begin{equation}
  V = R_m\,i + L_m\frac{\dd i}{\dd t} + K_e\,\omega_{\mathrm{rel}} .
\end{equation}
The electrical time constant $L_m/R_m$ of motors of this size is of the order
of \SI{1}{ms}, much shorter than the time scale of the robot
($1/\sqrt\beta = \SI{87}{ms}$, \cref{sec:lin}); the inductive term is
neglected, and
\begin{equation}
  i = \frac{V - K_e\,\omega_{\mathrm{rel}}}{R_m} .
\end{equation}

\paragraph{Torque.} The torque at the output shaft is proportional to the
current, $\tau_1 = K_t\,i$, hence
\begin{equation}
  \tau_1 = \frac{K_t}{R_m}\,V - \frac{K_t K_e}{R_m}\,\omega_{\mathrm{rel}} .
  \label{eq:motor-K}
\end{equation}

\paragraph{Relative speed.} The housing of the motor is fixed to the body,
which rotates with $\theta$; the motor drives the wheel relative to the body.
Using the rolling constraint $x = r\varphi$,
\begin{equation}
  \omega_{\mathrm{rel}} = \dot\varphi - \dot\theta = \frac{\dot x}{r} - \dot\theta ,
\end{equation}
the same combination that appears in the virtual work of the motor,
$\delta W = \tau(\delta\varphi - \delta\theta)$ (\cref{eq:Q}).

\paragraph{Constants from the datasheet.} At stall ($\omega_{\mathrm{rel}} = 0$,
$V = V_{\mathrm{nom}}$) the current is $i_s = V_{\mathrm{nom}}/R_m$ and the
torque is $\tau_s = K_t i_s$:
\begin{equation}
  R_m = \frac{V_{\mathrm{nom}}}{i_s} = \SI{9.23}{\ohm} ,
  \qquad K_t = \frac{\tau_s}{i_s} = \SI{0.128}{N.m/A} .
\end{equation}
At no load the torque, and so the current, is approximately zero (the
no-load current is 4.6\% of $i_s$ and is neglected), so
$V_{\mathrm{nom}} = K_e\,\omega_0$:
\begin{equation}
  K_e = \frac{V_{\mathrm{nom}}}{\omega_0} = \SI{0.409}{V.s/rad} .
\end{equation}
For a lossless motor and gearbox $K_t = K_e$ in SI units; here
$K_t/K_e = 0.31$, the difference collecting the losses of the gearbox and the
approximations of the datasheet. Substituting the three constants in
\cref{eq:motor-K},
\begin{equation}
  \frac{K_t}{R_m} = \frac{\tau_s}{V_{\mathrm{nom}}} ,
  \qquad \frac{K_t K_e}{R_m} = \frac{\tau_s}{i_s}\,\frac{V_{\mathrm{nom}}}{\omega_0}\,\frac{i_s}{V_{\mathrm{nom}}}
  = \frac{\tau_s}{\omega_0} ,
\end{equation}
and the torque of one motor is written with datasheet quantities only:
\begin{equation}
  \tau_1 = \tau_s\frac{V}{V_{\mathrm{nom}}} - \frac{\tau_s}{\omega_0}\left(\frac{\dot x}{r} - \dot\theta\right) .
  \label{eq:motor}
\end{equation}
The torque decreases linearly with the relative speed and is zero at
$\omega_{\mathrm{rel}} = \omega_0\,V/V_{\mathrm{nom}}$. With $V = 0$ (terminals
shorted by the driver) the second term brakes the wheel. The same voltage is
applied to both motors, so the total torque of \cref{sec:model} is
$\tau = 2\tau_1$.

\paragraph{Limits.} The voltage is limited by the battery,
$|V| \le V_b$, with $V_b$ between \SI{10.8}{V} (discharged) and \SI{12.6}{V}
(charged), \SI{11.1}{V} nominal. Below a threshold voltage the static friction
of the gearbox keeps the motor still; the threshold is estimated from the ratio
of no-load to stall current, $V_d = (0.06/1.3)\,V_{\mathrm{nom}} \approx \SI{0.5}{V}$,
and the voltage that drives the motor is
\begin{equation}
  V_{\mathrm{eff}} = \operatorname{sign}(V)\,\max\left(|V| - V_d,\ 0\right) ,
\end{equation}
which replaces $V$ in \cref{eq:motor}. At a given speed the available torque is
therefore bounded by the values of \cref{eq:motor} at $V = \pm V_b$. The maximum
speed of the robot on flat ground is $\omega_0 r = \SI{1.35}{m/s}$.

\subsection{Use in the simulation}

The controllers of \cref{sec:control} still compute a torque $\tau$. The
voltage that produces it at the current speed is obtained by inverting
\cref{eq:motor} and adding the dead-zone threshold,
\begin{equation}
  V = V_{\mathrm{nom}}\left(\frac{\tau}{2\tau_s} + \frac{\omega_{\mathrm{rel}}}{\omega_0}\right)
      + \operatorname{sign}(V)\,V_d ,
\end{equation}
limited to $\pm V_b$. The voltage is held constant during each step, while the
torque is recomputed from \cref{eq:motor} at every evaluation of the
dynamics inside the step, since it depends on the speed.

\subsection{Results}

From rest with $\theta_0 = 10^\circ$ and $V_b = \SI{11.1}{V}$, both controllers
balance the robot with the motor model. The PID uses at most \SI{5.6}{V} and
\SI{0.42}{A} per motor, the LQR \SI{6.2}{V} and \SI{0.61}{A}; the requested
torque stays at least \SI{0.14}{N.m} inside the available range.

The largest initial tilt from which each controller recovers (robot at rest,
upright within $1^\circ$ after \SI{5}{s}, tilts tested in steps of $2^\circ$)
is given in \cref{tab:roa} for several battery voltages. Without the motor
model, the same LQR recovers from up to $89^\circ$.

\begin{table}[h]
\centering
\caption{Largest recoverable initial tilt with the motor model.}
\label{tab:roa}
\begin{tabular}{@{}lcccc@{}}
\toprule
$V_b$ [V] & 10.8 & 11.1 & 12.0 & 12.6 \\
\midrule
PID & $34^\circ$ & $34^\circ$ & $38^\circ$ & $40^\circ$ \\
LQR & $34^\circ$ & $34^\circ$ & $38^\circ$ & $40^\circ$ \\
\bottomrule
\end{tabular}
\end{table}
